\input{../tex-inputs/latex-leseansicht-vorspann}

\section[Gustav Schwarzkopf an Arthur Schnitzler, 12. 9. 1899]{L04292 Gustav Schwarzkopf an Arthur Schnitzler, 12. 9. 1899}
\nopagebreak\mylabel{L04292v}
\rehead{ }\normalsize\beginnumbering\briefempfaengerindex{Schnitzler, Arthur@\textsc{Schnitzler, Arthur}!zzzSchwarzkopf, Gustav@\emph{von Gustav Schwarzkopf}!1899-09-123@{12. 9. 1899}|(be}
\toendnotes[C]{\smallbreak\pagebreak[2]}
\correspDesc{Versand  durch Gustav Schwarzkopf am 12. 9. 1899 in Brühl
\newline{}Erhalt  durch Arthur Schnitzler im Zeitraum [13. 9. 1899
                  – 16. 9. 1899?] in München}\toendnotes[C]{\smallbreak}
\Standort{CUL, Schnitzler, B 96a.}
\physDesc{Brief, 1 Blatt, 4 Seiten, 1462 Zeichen (Briefpapier mit gedrucktem Monogramm »IH\pwindex{Hiller, Ida 11.\,11.\,1858 Siklós – 3.\,8.\,1922 Wien@\textsc{Hiller, Ida} (11.\,11.\,1858 Siklós – 3.\,8.\,1922 Wien)|pw}«)
\newline{}Handschrift: schwarze Tinte, deutsche Kurrent}\toendnotes[C]{\smallbreak}
\pstart
           \raggedleft{}{\pb}12. 9. 99\pend
           \vspace{0.5em}
\pstart
           Lieber Arthur,{ }ſeit 3 Tagen bin ich in der Brühl\oindex{Brühl@\textbf{Brühl}, \emph{Tal}|pw} – auf allgemeines Verlangen – natürlich regnet es{ }ſeit
               diesen drei Tagen und hat 9°, weil \uline{ich} »köstliche
               Landluft athmen« und mich an »den{ }ſonnigen Reizen der Natur erfreuen will«. Ihr Brief
               wurde mir erſt heute nachgeſchickt, meine Antwort trifft hoffentlich noch zu rechter
               Zeit ein. – Ich glaube, Sie wollen und \uline{werden} die
               kleine Reiſe unternehmen. Warum auch nicht? {\pb}In Berlin\oindex{Berlin@\textbf{Berlin}, \emph{Hauptstadt}|pw} aber dürfte die Situation zitternd werden, we{\geminationn} auch vielleicht nicht i{\geminationm}er ganz gemütlich. Unter dieſen Umſtänden dürfte es wol Anfang October werden, bis
               Sie nach Wien\oindex{Wien@\textbf{Wien}, \emph{Verwaltungsgebiet}|pw} ko{\geminationm}en.\pend
           
\pstart
           Richard\pwindex{Beer-Hofmann, Richard 11.\,7.\,1866 Wien – 26.\,9.\,1945 New York City@\textsc{Beer-Hofmann, Richard} (11.\,7.\,1866 Wien – 26.\,9.\,1945 New York City), \emph{Schriftsteller}|pw} hat mir einen Brief geſchrieben,
               ausgerechnet \uline{einen} Tag vor{ }ſeiner Abreiſe von Seeboden\oindex{Seeboden am Millstättersee@\textbf{Seeboden am Millstättersee}|pw}. »Muſchelkinder\pwindex{\textcolor{red}{\textsuperscript{XXXX indx1}}!Muschelkinder. Schauspiel in vier Acten nach einem Roman von Guy de Maupassant@\strich\emph{Muschelkinder. Schauspiel in vier Acten nach einem Roman von Guy de Maupassant}|pw}« habe ich nicht geſehen, die letzten Nu{\geminationm}ern der »Fackel\pwindex{Fackel@\emph{Die Fackel}|pw}«
               nicht geleſen, und ein Drama habe ich {\pb}auch noch nicht bego{\geminationn}en. Allein aus unſerem kleinen
               Kreis heraus werden die deutſchen Bühnen für die nächſten Zeit{ }ſo reichlich verſorgt
               werden, daß{ }ſie mich wirklich nicht brauchen.\pend
           
\pstart
           Über die unſagbare Gemeinheit in der Sache \textsc{Dreyfus\pwindex{Dreyfus, Alfred 9.\,10.\,1859 Mulhouse – 12.\,7.\,1935 Paris@\textsc{Dreyfus, Alfred} (9.\,10.\,1859 Mulhouse – 12.\,7.\,1935 Paris), \emph{Militär}|pw}} werden wir uns noch gemeinſam mündlich genug ärgern,{ }ſchreiben ka{\geminationn} man darüber nicht.–\pend
           
\pstart
           Die »Zerſtreuung« wird hoffentlich intenſiv genug{ }ſein, um {\pb}Ihre \substVorne{}\textsuperscript{unei\textcolor{gray}{×}\-\textcolor{gray}{×}\-\textcolor{gray}{×}}\substDazwischen{}Sti{\geminationm}un\substHinten{}g zu verbeſſern und  Ihnen \introOben{}auch\introOben{} Ihr Stück\pwindex{Schnitzler, Arthur 15.\,5.\,1862 Wien – 21.\,10.\,1931 ebd.@\textsc{Schnitzler, Arthur} (15.\,5.\,1862 Wien – 21.\,10.\,1931 ebd.), \emph{Schriftsteller, Mediziner}!Schleier der Beatrice. Schauspiel in fünf Akten@\strich\emph{Der Schleier der Beatrice. Schauspiel in fünf Akten}|pwv} beſſer erſcheinen zu laſſen. Ich bin{ }ſchon
               neugierig.\pend
           
\pstart
           Von \textsc{Hiller’s\pwindex{Hiller, Ida 11.\,11.\,1858 Siklós – 3.\,8.\,1922 Wien@\textsc{Hiller, Ida} (11.\,11.\,1858 Siklós – 3.\,8.\,1922 Wien)|pw}\pwindex{Hiller, Max 17.\,9.\,1856 Osijek – 13.\,1.\,1941 Wien@\textsc{Hiller, Max} (17.\,9.\,1856 Osijek – 13.\,1.\,1941 Wien), \emph{Industrieller}|pw}}, in deren Wohnung\oindex{?? [Wohnung von Ida und Max Hiller, Brühl, 1899]@\textbf{?? [Wohnung von Ida und Max Hiller, Brühl, 1899]}, \emph{Wohnung}|pwv} ich{ }ſchreibe
               und deren Briefpapier ich benütze die beſten Grüße.\pend
           
\pstart
           We{\geminationn} Sie wollen, kö{\geminationn}en Sie
               den beiden Damen\pwindex{Elsinger, Marie *~28.\,2.\,1874 St. Pölten@\textsc{Elsinger, Marie} (*~28.\,2.\,1874 St. Pölten), \emph{Schauspielerin}|pwv}\pwindex{Glümer, Marie 3.\,7.\,1867 Wien – 16.\,11.\,1925 München@\textsc{Glümer, Marie} (3.\,7.\,1867 Wien – 16.\,11.\,1925 München), \emph{Schauspielerin}|pwv} von mir Empſehlungen
               beſtellen.\pend
           
\pstart
           Herzlichſt{\\[\baselineskip]}Ihr{\\[\baselineskip]}\spacefill\mbox{Gustav}\pend
           \leftskip=0em{}\selectlanguage{ngerman}\endnumbering\briefempfaengerindex{Schnitzler, Arthur@\textsc{Schnitzler, Arthur}!zzzSchwarzkopf, Gustav@\emph{von Gustav Schwarzkopf}!1899-09-123@{12. 9. 1899}|)be}\mylabel{L04292h}
\begin{anhang}
\end{anhang}\newcommand{\dateiname}{L04292}\newcommand{\titel}{Gustav Schwarzkopf an Arthur Schnitzler, 12. 9. 1899}\newcommand{\editorInnen}{Herausgegeben von Jahnke, SelmaMüller, Martin Anton}\input{../tex-inputs/latex-leseansicht-abspann}
